% !TEX root = main.tex
\section{はじめに：位相同期回路（PLL）はなぜ必要か}
\label{sec:intro}

安定した周波数の信号を作り出す回路は，一般に位相同期回路（Phase-Locked Loop, PLL）によって実現される．PLLは有線通信・無線通信・レーダなど，多くのシステムの根幹を支える部品であり，その性能がそのままシステム全体の性能の上限を決めてしまうことが少なくない．以下に代表的な3つの用途を挙げ，PLLの性能がどのように効いてくるかを整理する．

\begin{itemize}
  \item \textbf{有線通信（SerDes）：} 送信側のPLLがデータ送出のタイミングを作るクロックを生成し，受信側では再びPLLを含むCDR（Clock and Data Recovery）回路がデータを再サンプリングして0/1を判定する．このとき、サンプリングの判定が最も確実なのは、「出力クロックのエッジが、入力が遷移するちょうど中間にあるとき」である．ここで、PLLの出力クロックに時間方向の揺らぎ（ジッタ）が乗ると，受信側のサンプリング点がずれてビット誤り率（BER）が悪化する．
  \item \textbf{無線通信：} 受信信号を局部発振器（Local Oscillator, LO）の信号とミキシングして周波数を落とす際，LOの位相雑音が悪いと信号の品質を悪化させる．
\end{itemize}

このように，用途は異なっていても「PLLの生成する信号の位相雑音（ジッタ）がシステムの性能上限を決める」という構図は共通している．したがって，PLLを低雑音かつ低消費電力で実現することには大きな意味がある．

本報告書では，まずPLLの基本構成と，本研究の出発点となる問題（分数分周比PLLにおける量子化雑音の増幅）を第\ref{sec:pll-basics}節で整理し，低雑音化がなぜ重要かを第\ref{sec:motivation}節で述べる．続いて，位相雑音を抑制する既存手法と，本研究が対象とするフィードフォワード（FF）型構成を第\ref{sec:methods}節で紹介し，本研究の目的（第\ref{sec:objective}節）・提案する取り組み方（第\ref{sec:proposal}節）・これまでの進捗（第\ref{sec:progress}節）・今後の展望（第\ref{sec:future}節）を順に述べる．
